\documentclass[10pt,a4paper]{amsart}
\usepackage{amsmath,amssymb,amsthm,graphicx}
% Retained minimal preamble for the Stacks Project chapter more-algebra.tex (More on Algebra).
% The chapter file pinned-original-source.tex carries only \input{preamble} on its first line, so no
% preamble text lives inside the chapter itself.  The Stacks Project uses one single shared preamble for
% every chapter; the snapshot stages it identically as sources/stacks-chow/preamble.tex and
% sources/stacks-etale-cohomology/preamble.tex.  The block below is copied mechanically, line for line,
% from sources/stacks-chow/preamble.tex lines 14-18 (the Xy-pic setup) and 154-217 (the native theorem
% declarations together with the mathematical macros).  Nothing was added, removed or redefined; the
% wrapper class and the amsmath/amssymb/amsthm/graphicx base packages are supplied by the assembler.
% Chapter: more-algebra.tex -- More on Algebra.
% For commutative diagrams we use Xy-pic
\usepackage[all]{xy}

% We use 2cell for 2-commutative diagrams.
\xyoption{2cell}

% Theorem environments.
%
\theoremstyle{plain}
\newtheorem{theorem}[subsection]{Theorem}
\newtheorem{proposition}[subsection]{Proposition}
\newtheorem{lemma}[subsection]{Lemma}

\theoremstyle{definition}
\newtheorem{definition}[subsection]{Definition}
\newtheorem{example}[subsection]{Example}
\newtheorem{exercise}[subsection]{Exercise}
\newtheorem{situation}[subsection]{Situation}

\theoremstyle{remark}
\newtheorem{remark}[subsection]{Remark}
\newtheorem{remarks}[subsection]{Remarks}

\numberwithin{equation}{subsection}

% Macros
%
\def\lim{\mathop{\mathrm{lim}}\nolimits}
\def\colim{\mathop{\mathrm{colim}}\nolimits}
\def\Spec{\mathop{\mathrm{Spec}}}
\def\Hom{\mathop{\mathrm{Hom}}\nolimits}
\def\Ext{\mathop{\mathrm{Ext}}\nolimits}
\def\SheafHom{\mathop{\mathcal{H}\!\mathit{om}}\nolimits}
\def\SheafExt{\mathop{\mathcal{E}\!\mathit{xt}}\nolimits}
\def\Sch{\mathit{Sch}}
\def\Mor{\mathop{\mathrm{Mor}}\nolimits}
\def\Ob{\mathop{\mathrm{Ob}}\nolimits}
\def\Sh{\mathop{\mathit{Sh}}\nolimits}
\def\NL{\mathop{N\!L}\nolimits}
\def\CH{\mathop{\mathrm{CH}}\nolimits}
\def\proetale{{pro\text{-}\acute{e}tale}}
\def\etale{{\acute{e}tale}}
\def\QCoh{\mathit{QCoh}}
\def\Ker{\mathop{\mathrm{Ker}}}
\def\Im{\mathop{\mathrm{Im}}}
\def\Coker{\mathop{\mathrm{Coker}}}
\def\Coim{\mathop{\mathrm{Coim}}}

% Boxtimes
%
\DeclareMathSymbol{\boxtimes}{\mathbin}{AMSa}{"02}

%
% Macros for moduli stacks/spaces
%
\def\QCohstack{\mathcal{QC}\!\mathit{oh}}
\def\Cohstack{\mathcal{C}\!\mathit{oh}}
\def\Spacesstack{\mathcal{S}\!\mathit{paces}}
\def\Quotfunctor{\mathrm{Quot}}
\def\Hilbfunctor{\mathrm{Hilb}}
\def\Curvesstack{\mathcal{C}\!\mathit{urves}}
\def\Polarizedstack{\mathcal{P}\!\mathit{olarized}}
\def\Complexesstack{\mathcal{C}\!\mathit{omplexes}}
% \Pic is the operator that assigns to X its picard group, usage \Pic(X)
% \Picardstack_{X/B} denotes the Picard stack of X over B
% \Picardfunctor_{X/B} denotes the Picard functor of X over B
\def\Pic{\mathop{\mathrm{Pic}}\nolimits}
\def\Picardstack{\mathcal{P}\!\mathit{ic}}
\def\Picardfunctor{\mathrm{Pic}}
\def\Deformationcategory{\mathcal{D}\!\mathit{ef}}

% --- appended 2026-10-12 (Chief): the project own definitions of the reference/slogan/history
% display environments, copied mechanically from the shared Stacks preamble lines 8-12
% (\usepackage{verbatim} plus the three \newenvironment lines, which the chapter preamble keeps
% verbatim from the project preamble). Nothing added, removed or redefined: these environments
% carry no mathematical content and the project itself typesets them as comments. Without these
% three lines a counted statement whose own text uses them fails with Environment slogan undefined.
\usepackage{verbatim}
\newenvironment{reference}{\comment}{\endcomment}
\newenvironment{slogan}{\comment}{\endcomment}
\newenvironment{history}{\comment}{\endcomment}
\usepackage{cleveref}
\crefname{originallemma}{original lemma}{original lemmas}
\Crefname{originallemma}{Original lemma}{Original lemmas}
\makeatletter
\newlabel{lemma-restriction-derived-complete}{{0924}{source}{}{sourcepointer-1}{}}
\newlabel{lemma-restriction-derived-complete@cref}{{[originallemma][0][]0924}{[1][0][]source}{}{}{}}
\makeatother
\crefname{originalsection}{original section}{original sections}
\Crefname{originalsection}{Original section}{Original sections}
\makeatletter
\newlabel{example-derived-complete-modules}{{09AT}{source}{}{sourcepointer-2}{}}
\newlabel{example-derived-complete-modules@cref}{{[originalsection][0][]09AT}{[1][0][]source}{}{}{}}
\makeatother
\crefname{originallemma}{original lemma}{original lemmas}
\Crefname{originallemma}{Original lemma}{Original lemmas}
\makeatletter
\newlabel{lemma-consequence-baire-complete-module}{{0CQV}{source}{}{sourcepointer-3}{}}
\newlabel{lemma-consequence-baire-complete-module@cref}{{[originallemma][0][]0CQV}{[1][0][]source}{}{}{}}
\makeatother
\crefname{originallemma}{original lemma}{original lemmas}
\Crefname{originallemma}{Original lemma}{Original lemmas}
\makeatletter
\newlabel{lemma-open-mapping}{{0CQW}{source}{}{sourcepointer-4}{}}
\newlabel{lemma-open-mapping@cref}{{[originallemma][0][]0CQW}{[1][0][]source}{}{}{}}
\makeatother
\title{A derived complete module that is power torsion for a finitely generated ideal is annihilated by a power of that ideal}
\author{The Stacks Project Authors}
\date{stacks-project:more-algebra.tex:local-snapshot-20261012; GNU FDL 1.2 or later}
\begin{document}
\maketitle
\setcounter{subsection}{1}
\noindent\textit{Editorial note.} All delivered mathematical text is the author's own, copied from the pinned Stacks Project chapter file more-algebra.tex; no statement or proof step was rewritten, reordered or removed. The counted span is the author environment and its own proof environment. References inside the span are handed to the assembler as navigation stubs that carry the canonical tag of the cited tag only; whether each of them is load-bearing is recorded on the stub itself.\par\medskip
\noindent\textit{Allowed standard prerequisites (editorial).}\begin{itemize}
\item Derived completeness is preserved by restriction along a ring map carrying the ideal back, so it can be transferred to a polynomial algebra. invoked at proof lines 27024-27026. Lets the proof replace the ring by the polynomial ring in one variable and assume that the generator is a nonzerodivisor.
\item Every module sits in a two-term sequence whose terms are adically complete modules with no torsion along the generator. invoked at proof lines 27031-27034. Provides the surjection from the complete module L whose kernel K is also complete and torsion free, which carries the whole argument.
\item In a complete topological module, a countable union of closed subsets covering the whole module has a member with interior. invoked at proof lines 27045-27047. Makes one of the closed pieces N\_n open in L, the key step that lets the argument climb from K to L.
\item An open continuous surjection of complete topological modules has a controlled image, so a power of the generator lands inside the image. invoked at proof lines 27055-27058. Produces the exponent t with the t-th power of f contained in the image of u, which is the stated annihilation.
\end{itemize}\medskip
\ifdefined\hypertarget\hypertarget{sourcepointer-1}{}\fi
\noindent\textit{Source reference.} Original lemma 0924: sources/stacks-more-algebra/pinned-original-source.tex, source lines 26616--26616. Invoked at proof lines 27024-27026 to see that M is derived complete as a module over the polynomial ring with respect to the principal ideal generated by x; the use of that transfer, namely the reduction to the case of a nonzerodivisor, is inside the counted span.\par\medskip
\ifdefined\hypertarget\hypertarget{sourcepointer-2}{}\fi
\noindent\textit{Source reference.} Original section 09AT: sources/stacks-more-algebra/pinned-original-source.tex, source lines 26844--26844. Cited at proof lines 27031-27034 for the existence of the sequence with K and L adically complete and torsion free, and the footnote re-proves the only part that is used; every later step of the argument is inside the counted span.\par\medskip
\ifdefined\hypertarget\hypertarget{sourcepointer-3}{}\fi
\noindent\textit{Source reference.} Original lemma 0CQV: sources/stacks-more-algebra/pinned-original-source.tex, source lines 9078--9078. Applied at proof lines 27045-27047 to the covering of L by the closures N\_n to obtain an index n with N\_n open; the construction of the increasing pieces L\_n and the containments verified around it are inside the counted span.\par\medskip
\ifdefined\hypertarget\hypertarget{sourcepointer-4}{}\fi
\noindent\textit{Source reference.} Original lemma 0CQW: sources/stacks-more-algebra/pinned-original-source.tex, source lines 9090--9090. Used at proof lines 27055-27058 to conclude that the map u is open from the openness of the closures of the images; establishing that openness by the continuity and openness of multiplication by a power of f is inside the counted span.\par\medskip
\begin{lemma}[Bhatt]
\label{lemma-torsion-and-derived-complete}
Let $I$ be a finitely generated ideal in a ring $A$.
Let $M$ be a derived complete $A$-module. If $M$ is
an $I$-power torsion module, then $I^nM = 0$ for some $n$.
\end{lemma}

\begin{proof}
Say $I = (f_1, \ldots, f_r)$. It suffices to show
that for each $i$ there is an $n_i$ such that $f_i^{n_i}M = 0$.
Hence we may assume that $I = (f)$ is a principal ideal.
Let $B = \mathbf{Z}[x] \to A$ be the ring map sending $x$ to $f$.
By Lemma \ref{lemma-restriction-derived-complete}
we see that $M$ is derived complete as a $B$-module
with respect to the ideal $(x)$. After replacing $A$ by
$B$, we may assume that $f$ is a nonzerodivisor in $A$.

\medskip\noindent
Assume $I = (f)$ with $f \in A$ a nonzerodivisor.
According to Example \ref{example-derived-complete-modules}
there exists a short exact sequence
$$
0 \to K \xrightarrow{u} L \to M \to 0
$$
where $K$ and $L$ are $I$-adically complete $A$-modules
whose $f$-torsion is zero\footnote{For the proof it is enough
to show that there exists a sequence $K \xrightarrow{u} L \to M \to 0$
where $K$ and $L$ are $I$-adically complete $A$-modules. This can
be shown by choosing a presentation $F_1 \to F_0 \to M \to 0$
with $F_i$ free and then setting $K$ and $L$ equal to the
$f$-adic completions of $F_1$ and $F_0$. Namely, as $f$
is a nonzerodivisor these completions will be the
derived completions and the sequence will remain exact.}.
Consider $K$ and $L$ as
topological modules with the $I$-adic topology. Then $u$ is continuous.
Let
$$
L_n = \{x \in L \mid f^n x \in u(K)\}
$$
Since $M$ is $f$-power torsion we see that
$L = \bigcup L_n$. Let $N_n$ be the closure of $L_n$ in $L$.
By Lemma \ref{lemma-consequence-baire-complete-module}
we see that $N_n$ is open in $L$ for some $n$. Fix such an $n$.
Since $f^{n + m} : L \to L$ is a continuous open map, and since
$f^{n + m} L_n \subset u(f^m K)$ we conclude that
the closure of $u(f^mK)$ is open for all $m \geq 1$.
Thus by Lemma \ref{lemma-open-mapping}
we conclude that $u$ is open. Hence $f^tL \subset \Im(u)$
for some $t$ and we conclude that $f^t$ annihilates $M$
as desired.
\end{proof}

\end{document}
